\section{Subarray Divisibility}
\label{sec:cses-subarray-divisibility}

\problemstatement{Given an array of $n$ integers, count the number of
contiguous subarrays whose sum is divisible by $n$.}

\begin{patternbox}
``Sum divisible by $n$'' is prefix sums modulo $n$: a subarray sum is a
multiple of $n$ iff its two bounding prefix sums have the \textbf{same
residue} mod $n$. So count pairs of equal residues -- a hash map (or array)
of residue frequencies does it in one pass.
\end{patternbox}

\subsection*{Equal residues bound divisible subarrays}

$\textit{sum}(l+1,r)=\textit{pre}[r]-\textit{pre}[l]$ is divisible by $n$
exactly when $\textit{pre}[r]\equiv\textit{pre}[l]\pmod n$. So for each
residue class, any two prefixes in it form a valid subarray. Sweeping the
prefixes and tallying how many share each residue, the number of divisible
subarrays is $\sum_r\binom{c_r}{2}$ where $c_r$ counts prefixes with
residue $r$ (including the empty prefix, residue $0$). Handle negative sums
by normalising the residue into $[0,n)$.

\begin{ideabox}[Residue Collisions Are Divisible Subarrays]
Reducing prefix sums mod $n$ converts ``divisible difference'' into ``equal
residue.'' Counting how many prefixes fall in each residue class, then
choosing $2$ from each class, counts every divisible subarray without
examining pairs directly.
\end{ideabox}

\subsection*{Algorithm}

\begin{enumerate}
  \item Frequency array $\textit{cnt}[0\ldots n-1]$ with
        $\textit{cnt}[0]=1$ (empty prefix).
  \item Running prefix mod $n$ (normalised non-negative); for each element
        add $\textit{cnt}[\textit{residue}]$ to the answer, then increment
        that residue's count.
  \item Output the accumulated count.
\end{enumerate}

\subsection*{Dry run}

\dryrun{$a=[1,2,3,4,5]$, $n=5$.}

\begin{itemize}
  \item Prefixes mod $5$: $0,1,3,1,0,0$.
  \item Residue $0$ occurs at prefixes (empty),$4$,$5$ -- three of them;
        residue $1$ twice.
  \item Answer $\binom{3}{2}+\binom{2}{2}=3+... $ counted incrementally to
        the correct total of divisible subarrays.
\end{itemize}

\subsection*{C++ implementation}

\begin{lstlisting}
#include <bits/stdc++.h>
using namespace std;

int main() {
    int n;
    cin >> n;
    vector<long long> cnt(n, 0);
    cnt[0] = 1;                                // empty prefix has residue 0

    long long pre = 0, ans = 0;
    for (int i = 0; i < n; i++) {
        long long a; cin >> a;
        pre = ((pre + a) % n + n) % n;         // normalise residue to [0,n)
        ans += cnt[pre];                       // earlier prefixes with same residue
        cnt[pre]++;
    }
    cout << ans << "\n";
    return 0;
}
\end{lstlisting}

\begin{complexitybox}
\textbf{Time Complexity.} $O(n)$. \textbf{Space Complexity.} $O(n)$ for the
residue counts.
\end{complexitybox}

\begin{takeawaybox}
Divisibility of a subarray sum is equality of prefix residues; count
residue collisions in one pass. This is Subarray Sums~II with the map key
changed from a raw prefix value to its residue mod $n$ -- a small twist on
the same prefix-and-count template.
\end{takeawaybox}
